% Chapter 04 overview
\section{ওয়েব ডিজাইন পরিচিতি এবং HTML — পরিদর্শন}

\begin{learningbox}
এই অধ্যায়ে আমরা ওয়েব ডিজাইনের মূল ধারণা, HTML-এর ভূমিকা, ওয়েবসাইটের কাঠামো ও পেজ ডিজাইন প্রক্রিয়া আলোচনা করব। অধ্যায় শেষে শিক্ষার্থী ওয়েব পেজ তৈরি ও publish-এর মৌলিক ধাপ বুঝতে পারবে।
\end{learningbox}

\subsection{চাপ ও লক্ষ্য}
এই অধ্যায়টি ওয়েব ডেভেলপমেন্টের প্রথম ধাপ হিসেবে কাজ করবে। এখানে HTML দিয়ে কিভাবে একটি static page তৈরি হয়, কিভাবে তা structure করা উচিত এবং কোন tag কোন কাজে ব্যবহৃত হবে—এসব শিখানো হবে। ছোট উদাহরণ, layout নীতিমালা এবং practical task-ও থাকবে।

\subsection{কী শেখা হবে}
\begin{itemize}
\item ওয়েব পেজ ও ওয়েবসাইটের মৌলিক ধারণা।
\item HTML document structure ও সাধারণ tag-এর ব্যবহার।
\item Text formatting, links, images, tables ও form-এর মাধ্যমে কনটেন্ট উপস্থাপন।
\item ওয়েব পেজ design-এর workflow এবং publish করার স্টেপ।
\end{itemize}

\begin{practicebox}
\begin{enumerate}
\item ছোট একটি HTML skeleton লিখে browser-এ দেখাও (offline)।
\item একটি school homepage-এর wireframe আঁকো এবং কোন tag কোথায় লাগবে তা সংক্ষেপে লেখো।
\end{enumerate}
\end{practicebox}
